\chapter{Simulation}
\label{ch:simulation}

This chapter describes the simulation chain of ATTPCROOTv2, from the generation of a beam
particle to the digitized pad traces that the reconstruction reads. The chain has two stages that
run as separate jobs:

\begin{enumerate}
  \item \emph{Transport.} A \code{FairRunSim} job generates primaries with one or more generators,
    transports them through the detector geometry with a Virtual Monte Carlo (VMC) engine
    (normally Geant4) and writes the energy deposits in the gas as \code{AtMCPoint} objects in the
    branch \code{AtTpcPoint}. MC tracks go to the branch \code{MCTrack}.
  \item \emph{Digitization.} A \code{FairRunAna} job reads that file and runs
    \code{AtClusterizeTask} (ionization electrons, drift, diffusion) and \code{AtPulseTask}
    (avalanche gain, electronics response, noise) to produce an \code{AtRawEvent} in the same
    format as unpacked experimental data. From that point the reconstruction tasks (PSA, pattern
    recognition, fitting) are the same as for data.
\end{enumerate}

A lighter path, \code{AtSimpleSimulation}, replaces the VMC transport with straight-line stepping
through an energy-loss model and writes \code{AtMCPoint} objects directly
(Section~\ref{sec:sim:simple}). Everything downstream of \code{AtMCPoint} is shared.

The source directories involved are \code{AtGenerators/} (event generators),
\code{AtSimulationData/} (\code{AtMCPoint}, \code{AtMCTrack}, \code{AtStack},
\code{AtVertexPropagator}), \code{AtDetectors/AtTpc/} (the sensitive detector),
\code{AtDigitization/} (clusterization, pulse generation, trigger) and \code{AtTools/} (electronic
response functions and space-charge models). The worked example in
Section~\ref{sec:sim:example} is the 16C(d,t)15C simulation in
\url{macro/Simulation/ATTPC/16C_dt/}.

%=====================================================================================
\section{Transport set-up in a simulation macro}
\label{sec:sim:transport}

\subsection{The FairRunSim skeleton}

Every transport macro in \url{macro/Simulation/ATTPC/} follows the same pattern. The excerpt
below is from \url{macro/Simulation/ATTPC/16C_dt/C16_dt_sim.C}:

\begin{verbatim}
FairRunSim *run = new FairRunSim();
run->SetName(mcEngine);                 // "TGeant4"
run->SetOutputFile(outFile);
FairRuntimeDb *rtdb = run->GetRuntimeDb();

run->SetMaterials("media.geo");

FairModule *cave = new AtCave("CAVE");
cave->SetGeometryFileName("cave.geo");
run->AddModule(cave);

FairDetector *ATTPC = new AtTpc("ATTPC", kTRUE);
ATTPC->SetGeometryFileName("ATTPC_D300torr_v2.root");
run->AddModule(ATTPC);

AtConstField *fMagField = new AtConstField();
fMagField->SetField(0., 0., bFieldkG);                 // kG
fMagField->SetFieldRegion(-50, 50, -50, 50, -10, 230); // cm
run->SetField(fMagField);
\end{verbatim}

The pieces are:

\begin{itemize}
  \item \textbf{Engine.} \code{SetName("TGeant4")} selects Geant4. The engine is configured by
    \url{gconfig/g4Config.C}, which builds a \code{TG4RunConfiguration} with the physics list
    \code{QGSP\_FTFP\_BERT} and the special processes
    \code{stepLimiter+specialCuts+specialControls+stackPopper}, creates the stack
    (\code{new AtStack(1000)}), and calls \code{SetMaxNStep(5E3)}. It then executes
    \url{gconfig/g4config.in}, which sets, among others, \code{/mcPhysics/rangeCuts 0.01 mm},
    \code{/mcPhysics/rangeCutForElectron 100 km} (no delta electrons are produced as separate
    tracks) and \code{/mcDet/setMaxStepInLowDensityMaterials 1}.
  \item \textbf{Media.} \code{SetMaterials("media.geo")} loads \url{geometry/media.geo}. The gas
    medium is referenced by name inside the geometry macro (for example
    \code{TargetD2\_300}, D$_2$ at 300~torr, density $6.5643\times10^{-5}$~g/cm$^3$).
  \item \textbf{Cave.} \code{AtCave} with \code{cave.geo} is the world volume.
  \item \textbf{Detector.} \code{AtTpc} is the sensitive detector. Its geometry is a ROOT file
    produced by a macro in \url{geometry/} (here \url{geometry/ATTPC_D300torr_v2.C}). In that
    macro the volume \code{drift\_volume} is a tube of diameter 50~cm and length 100~cm placed
    from $z=0$ to $z=100$~cm, with the entrance window at $z=0$.
  \item \textbf{Field.} \code{AtConstField} takes the field components in kG and the region in
    cm. A value of $-28.5$~kG is 2.85~T along $-z$.
\end{itemize}

After the generators are attached, the macro calls \code{run->Init()}, writes the parameter
file with a \code{FairParRootFileIo}, and runs \code{run->Run(nEvents)}.

\subsection{The sensitive detector: AtTpc}

A volume is sensitive when \code{AtTpc::CheckIfSensitive} finds \code{drift\_volume},
\code{window} or \code{cell} in its name. For every Geant4 step in such a volume,
\code{AtTpc::ProcessHits} records the step and adds one \code{AtMCPoint} to the
\code{TClonesArray} registered as \code{AtTpcPoint}. Each point stores:

\begin{itemize}
  \item the track ID (\code{GetTrackID}), the volume name (\code{GetVolName}) and copy number;
  \item position (cm) and momentum (GeV/c) at the step, time (ns), track length (cm);
  \item the energy deposited in the step (\code{GetEnergyLoss}, GeV);
  \item the mass and atomic numbers decoded from the PDG code (\code{GetMassNum},
    \code{GetAtomicNum});
  \item the initial energy and lab angle of the track as set by the reaction generator
    (\code{GetEIni} in MeV, \code{GetAIni} in degrees). For \code{fTrackID == 0} both are written as
    zero; for other tracks they are read from \code{AtVertexPropagator::GetTrackEnergy(trackID)} and
    \code{GetTrackAngle(trackID)}.
\end{itemize}

\code{AtTpc::SetDriftMaxStep(cm)} optionally limits the Geant4 step length in the gas medium. The
comments in \code{AtTpc::Initialize} state that with TGeant4 the \code{Gstpar} route is ignored and
the UI command route is rejected when issued from \code{Initialize}; the setter is best-effort and
does nothing unless called.

\subsection{AtStack}

\code{AtStack} (\url{AtSimulationData/AtStack.h}) is the VMC stack created in
\code{g4Config.C}. It counts MC points per track and detector (\code{AddPoint}, called by
\code{AtTpc::ProcessHits}) and, at the end of the event, selects which particles are written to the
\code{MCTrack} branch. Primaries are always kept. Secondaries are dropped if
\code{StoreSecondaries(kFALSE)} was called (default \code{kTRUE}), if they have fewer points than
\code{SetMinPoints} (default 1), or if their kinetic energy is below \code{SetEnergyCut} (GeV,
default 0). \code{StoreMothers} (default \code{kTRUE}) keeps the mothers of every selected track.
In the code these cuts act in \code{AtStack::SelectTracks}, i.e.\ on what is stored, not on what
Geant4 transports.

\subsection{Beam events, reaction events and AtVertexPropagator}
\label{sec:sim:vertex}

The AT-TPC simulation represents one beam-induced reaction as \emph{two consecutive FairRoot
events}:

\begin{itemize}
  \item a \textbf{beam event} (even entries 0, 2, 4, \ldots), in which only the beam ion is
    transported and is stopped at the reaction point;
  \item a \textbf{reaction event} (odd entries 1, 3, 5, \ldots), in which only the reaction products
    are transported, starting from that point.
\end{itemize}

\code{run->Run(2000)} therefore produces 1000 reactions at most. Note that the reaction event does
not contain the beam track and the beam event does not contain the products.

The state is carried between the two events by the singleton \code{AtVertexPropagator}
(\url{AtSimulationData/AtVertexPropagator.h}). It holds the beam/reaction flag
(\code{IsBeamEvent}, \code{IsReactionEvent}, toggled by \code{EndEvent()}), the vertex position and
the beam momentum and kinetic energy at the vertex (\code{SetVertex}, \code{GetVx/Vy/Vz},
\code{GetPx/Py/Pz}, \code{GetEnergy}), the random energy-loss threshold (\code{SetRndELoss},
\code{GetRndELoss}), the beam mass (\code{SetBeamMass}), the per-track generated energies and angles
(\code{SetTrackEnergy}, \code{SetTrackAngle}), and the momenta and excitation energies of the
reaction products for chained decay generators (\code{SetScatterP}, \code{SetRecoilP},
\code{SetScatterEx}, \code{SetRecoilEx}).

\paragraph{Where the reaction happens.} The mechanism involves three classes:

\begin{enumerate}
  \item On a beam event, \code{AtTPCIonGenerator::ReadEvent} samples a threshold
    $E_r$ uniformly in $[0,\ \mathrm{maxELoss}]$ (MeV) and stores it with
    \code{SetRndELoss(Er)}. It then adds the beam track.
  \item During stepping, \code{AtTpc} accumulates the energy deposited by the current track
    (\code{fELossAcc}, reset when the track enters a volume). \code{AtTpc::reactionOccursHere()}
    returns true when
\begin{verbatim}
bool atEnergyLoss =
   fELossAcc * 1000 > AtVertexPropagator::Instance()->GetRndELoss();
bool isPrimaryBeam = fTrackID == 0;
bool isInRightVolume =
   fVolName.Contains("drift_volume") || fVolName.Contains("cell");
\end{verbatim}
  \item \code{AtTpc::startReactionEvent()} then calls \code{gMC->StopTrack()}, resets the
    propagator, and stores the stop position, the beam momentum and the kinetic energy
    $E = (E_\mathrm{tot} - m_\mathrm{beam}\times 0.93149401)\times1000$~MeV with
    \code{SetVertex}. In the following reaction event the reaction generator reads these values.
\end{enumerate}

If the beam (track 0) leaves the drift volume without reacting, \code{AtTpc} calls
\code{ResetVertex()}, which sets the vertex energy to zero, and \code{AtTPC2Body} generates no
products in the reaction event (it requires \code{fBeamEnergy > 0}).

Because the accumulated loss in \code{fELossAcc} counts only the deposit inside the sensitive
volume, the vertex position follows from the stopping power of the gas: the threshold is uniform in
energy loss, which corresponds to a vertex distribution close to uniform in depth only while the
stopping power is close to constant over the chamber. Too large a \code{maxELoss} lets most beams
punch through without reacting; too small a value concentrates all vertices near the window. The
comment in \code{C16\_dt\_sim.C} states this, and recommends checking on the output that the vertex
distribution is flat along the drift length.

\paragraph{The residual and track ID 0.} In the reaction event no beam track is added, so the first
product pushed to the stack is also \code{fTrackID == 0}. \code{AtTpc} cannot distinguish it from the
beam. The guard is in \code{AtTPCIonGenerator::ReadEvent}: on reaction events it sets the threshold
to \code{std::numeric\_limits<double>::max()}, so \code{reactionOccursHere()} never fires on the
residual. The comment in \url{AtGenerators/AtTPCIonGenerator.cxx} documents the earlier behaviour
(commit \code{ac1b8fb1}, ``Disarm the reaction threshold on reaction events''): the threshold kept
its beam-event value, the heavy residual was stopped once it had deposited that same energy, and
on a 46Ar(3He,d) test only 46.8\% of the 47K residuals reached the end of the drift volume. The
light ejectile (track 1) was not affected. Simulations produced before that commit carry truncated
heavy residuals. The comment in \code{AtTpc::reactionOccursHere} also warns not to guard the test
with \code{AtVertexPropagator::IsBeamEvent()}: the flag is toggled during generation, so at stepping
time its sense is inverted.

%=====================================================================================
\section{Event generators}
\label{sec:sim:generators}

Generators are added to a \code{FairPrimaryGenerator} in the order they should run:

\begin{verbatim}
FairPrimaryGenerator *primGen = new FairPrimaryGenerator();
primGen->AddGenerator(ionGen);   // beam: AtTPCIonGenerator
primGen->AddGenerator(TwoBody);  // reaction: AtTPC2Body
run->SetGenerator(primGen);
\end{verbatim}

\subsection{AtReactionGenerator: the base class}

\code{AtReactionGenerator} (\url{AtGenerators/AtReactionGenerator.h}) derives from
\code{FairGenerator}. Subclasses implement \code{GenerateReaction(FairPrimaryGenerator*)}.
\code{ReadEvent} is declared \code{final}:

\begin{verbatim}
Bool_t AtReactionGenerator::ReadEvent(FairPrimaryGenerator *primGen)
{
   bool isBeamEvent = AtVertexPropagator::Instance()->IsBeamEvent();
   if (kIsFinalGen)
      // End the event if this is the final generator
      AtVertexPropagator::Instance()->EndEvent();

   if (!isBeamEvent || kAlwaysRun)
      return GenerateReaction(primGen);
   return true;
}
\end{verbatim}

A reaction generator therefore does nothing on beam events (unless \code{SetAlwaysRun(true)}), and
the last generator of the chain toggles the beam/reaction flag. \code{SetSequentialDecay(true)}
marks a generator as \emph{not} the last one (\code{kIsFinalGen = false}), so that a following
decay generator (for example \code{AtTPCIonDecay}) can read the products from
\code{AtVertexPropagator} and end the event. Exactly one generator in the chain must have
\code{kIsFinalGen} set, otherwise the alternation breaks; the header comment of
\code{AtVertexPropagator::EndEvent} states that it must be called once per event and that the user
is responsible for setting up the chain correctly.

\subsection{AtTPCIonGenerator: the beam}

\code{AtTPCIonGenerator} is a \code{FairGenerator} (not an \code{AtReactionGenerator}). The
constructor used by the AT-TPC macros is

\begin{verbatim}
AtTPCIonGenerator(const char *name, Int_t z, Int_t a, Int_t q, Int_t mult,
                  Double_t px, Double_t py, Double_t pz,
                  Double_t Ex, Double_t m, Double_t ener, Double_t eLoss = -1);
\end{verbatim}

\begin{itemize}
  \item \code{px, py, pz} are \emph{per nucleon} in GeV/c; the constructor multiplies them by
    \code{a}. A macro that knows the total momentum passes \code{pTot / a}.
  \item \code{m} is passed to \code{FairIon} as the ion mass. \code{FairIon} documents this
    argument in GeV; all AT-TPC macros pass it in atomic mass units, and the AT-TPC code that uses
    the stored value (\code{AtTpc::startReactionEvent}, \code{AtTPC2Body}) multiplies it by
    0.93149401, i.e.\ treats it as amu. See the pitfall list in Section~\ref{sec:sim:pitfalls}.
  \item \code{ener} is stored with \code{AtVertexPropagator::SetBeamNomE}. It does not set the
    beam energy; only the momentum does. The macro \url{macro/Simulation/ATTPC/17C_dp/C17_dp_sim.C}
    notes that \code{GetBeamNomE()} has no consumer in the tree.
  \item \code{eLoss} is \code{maxELoss}, the upper edge of the uniform threshold described in
    Section~\ref{sec:sim:vertex}. If it is negative (the default) the constructor uses \code{ener}
    instead, so omitting it makes the nominal beam energy the reaction threshold.
\end{itemize}

Setters: \code{SetSpotRadius(r, z, offx, offy)} gives the beam start: $x,y$ uniform in radius
$[0,r]$ and in azimuth around \code{(offx, offy)}, at fixed $z$ (all cm).
\code{SetDoReaction(false)} sets the threshold to the largest double so the beam never reacts and
adds the beam track on both beam and reaction events. \code{SetExcitationEnergy}, \code{SetMass} and
\code{SetCharge} modify the ion.

Two variants override \code{SetVertexCoordinates()}:
\begin{itemize}
  \item \code{AtTPCIonGeneratorGaussian}: \code{SetBeamLimits(r, z, theta)} and
    \code{SetBeamOrigin(x, y)}; the radius is a Gaussian of width $r/3$ clamped to $[0,r]$ and the
    direction is tilted by a uniform angle in $[0,\theta]$. As written, its constructor forwards
    \code{eLoss} into the \code{ener} slot of the base constructor, and \code{SetVertexCoordinates}
    multiplies the member \code{fPz} by $\cos\theta$ on every call, so with a non-zero
    \code{theta} the beam momentum decreases from event to event. Check the beam energy in the MC
    truth before using it with \code{theta > 0}.
  \item \code{AtTPCIonGeneratorS800}: beam emittance and momentum acceptance from S800 angular
    distributions (\code{SetBeamEmittance}, two histogram file names in the constructor).
\end{itemize}

\subsection{AtTPC2Body: two-body reactions}

\code{AtTPC2Body} (\url{AtGenerators/AtTPC2Body.h}) generates a binary reaction at the vertex left
by the beam. Constructor:

\begin{verbatim}
AtTPC2Body(const char *name, std::vector<Int_t> *z, std::vector<Int_t> *a,
           std::vector<Int_t> *q, Int_t mult, std::vector<Double_t> *px,
           std::vector<Double_t> *py, std::vector<Double_t> *pz,
           std::vector<Double_t> *mass, std::vector<Double_t> *Ex,
           Double_t ResEner, Double_t MinCMSAng, Double_t MaxCMSAng);
\end{verbatim}

\paragraph{The multiplicity-4 vector convention.} All vectors have four entries and \code{mult}
must be 4; \code{GenerateReaction} indexes entries 0--3 unconditionally.

\begin{center}
\begin{tabular}{p{1.3cm}p{3.6cm}p{8.5cm}}
\toprule
Index & Role (16C(d,t) macro) & Use in the code \\
\midrule
0 & beam (16C) & mass $m_1$ in the kinematics; momentum taken from \code{AtVertexPropagator} \\
1 & target (d) & mass $m_2$ \\
2 & heavy residual (15C) & mass $m_3$; gets $\theta_\mathrm{cm}$ sampled in the requested range;
  first track added to the stack (track ID 0 in the reaction event) \\
3 & light ejectile (t) & mass $m_4$; gets $\pi-\theta_\mathrm{cm}$; second track (track ID 1) \\
\bottomrule
\end{tabular}
\end{center}

Masses are in amu (the code computes $m\times 1000\times 0.93149401 + E_x$ in MeV) and
\code{Ex} in MeV. Only entries 2 and 3 are pushed to the stack. Entries with $A\neq1$ are created as
\code{FairIon}; $A=1,Z=1$ is a proton (PDG 2212) and $A=1,Z=0$ a neutron (PDG 2112).

\paragraph{Which particle the angular range refers to.} $\theta_\mathrm{cm}$ is sampled uniformly in
$\cos\theta$ between \code{MinCMSAng} and \code{MaxCMSAng} (degrees), i.e.\ isotropically in the
centre of mass, and is assigned to entry 2. The particle in entry 3 is emitted at
$\theta_\mathrm{cm,4} = 180^\circ - \theta_\mathrm{cm}$. In a macro that lists the heavy residual in
entry 2 and the light ejectile in entry 3 (all AT-TPC macros quoted here), the angular range is the
\emph{residual's} centre-of-mass angle. The header of
\url{macro/Simulation/ATTPC/16C_dt/run_states_dt.sh} uses this explicitly: it generates 2--70$^\circ$,
which covers ejectile centre-of-mass angles of 110--178$^\circ$.

The azimuth is uniform; entry 3 is emitted at $\phi+\pi$. Both directions are rotated from the beam
frame into the lab frame with \code{AtEulerTransformation} using the beam direction at the vertex.
If the centre-of-mass energy is below threshold, \code{SetValidKine(kFALSE)} is set and the products
get zero momentum. For each event the generator stores the lab energies and angles with
\code{SetTrackEnergy/SetTrackAngle}, key 0 for entry 2 and key 1 for entry 3; these become
\code{EIni}/\code{AIni} in the \code{AtMCPoint}s, matched by track ID. It also stores
\code{SetScatterP} (entry 2) and \code{SetRecoilP} (entry 3) for decay generators.

\paragraph{Options.} \code{SetFixedTargetPosition(vx, vy, vz)} switches to a fixed-target mode
(e.g.\ HELIOS-like set-ups): the vertex is fixed, \code{kAlwaysRun} is set, and the beam kinetic
energy is \code{ResEner}. In the active-target mode \code{ResEner} is unused.
\code{SetFixedBeamMomentum(px,py,pz)} replaces the propagated beam momentum.
\code{SetSequentialDecay(true)} removes entry 2 from the stack (it is expected to decay through the
next generator).

\subsection{Other generators}

\begin{center}
\begin{tabular}{p{4.3cm}p{10.4cm}}
\toprule
Class & What it does (from the header and source) \\
\midrule
\code{AtTPCIonDecay} & \code{AtReactionGenerator}; phase-space decay (\code{TGenPhaseSpace})
  of an excited product into one of several channels with given separation energies; reads the
  product momentum and excitation from \code{AtVertexPropagator} when used after a reaction
  generator with \code{SetSequentialDecay(true)}. \\
\code{AtTPCReactionDecay} & \code{AtReactionGenerator}; scatter and recoil described by a
  \code{DecayIon} struct (Z, A, q, masses, multiplicity, $E_x$, separation energy), each
  decayed by phase space. \\
\code{AtTPCIonPhaseSpace} & \code{FairGenerator}; N-body phase space of the beam--target system. \\
\code{AtTPC\_d2He} & \code{AtReactionGenerator}; the (d,$^2$He) reaction with a cross-section table
  passed as three vectors plus the table length \code{N\_data}. \\
\code{AtTPC\_Background} & \code{AtReactionGenerator}; background from deuteron elastic breakup
  (written by J.~C.~Zamora); the source treats the deuteron as unbound by up to 1~MeV and uses an
  internal two-body helper (\code{TwoB}). \\
\code{AtTPCXSReader} & \code{AtReactionGenerator}; energy and angle drawn with \code{GetRandom2}
  from a \code{TH2F} filled from a fixed-format text file ($31\times18$ values, PACE cross sections)
  in \url{AtGenerators/}, for example \url{xs_22Mgp_fusionEvaporation.txt}. \\
\code{AtTPCXSManager} & singleton holding a \code{TH2F} excitation function
  (\code{SetExcitationFunction(file)}). \code{AtTPCXSReader} does not use it; the only user in the
  tree is \url{resources/cross_sections/test_XSManager.C}. \\
\code{AtTPCFission\-Generator}, \code{AtTPCFission\-GeneratorV2}, \code{AtTPCFission\-GeneratorV3} &
  fission fragments read from a ROOT tree of pre-computed events (branches \code{Evnt},
  \code{Ntrack}, \code{Aout}, \code{Zout}, \code{fOutPx}, \ldots\ in V1 and V2). V3 takes an ion
  list and a fragment file, and boosts the fragments with the beam at the vertex to the lab. \\
\code{AtTPC20MgDecay}, \code{AtTPC20MgDecay\_pag} & $^{20}$Mg decay: isotropic protons, alphas,
  gammas and betas at hard-coded energies; point vertex (\code{SetXYZ}) or box vertex
  (\code{SetBoxXYZ}), \code{ShowOnlyAlphaProtonBranch}, decay chain set with
  \code{SetDecayChainPoint(E, BR)}. \\
\code{AtTPCGamma\-Dummy\-Generator} & \code{FairGenerator} based on an R3B test generator; particles
  of a given PDG with momentum, angle and vertex ranges (\code{SetPRange}, \code{SetThetaRange},
  \code{SetXYZ}, \ldots). \\
\code{AtPUMAGenerator} & \code{FairGenerator}; $\bar p$ + $^3$He at rest decayed by phase space to
  K$^+$K$^+\pi^-$ + hexaquark (PDG set with \code{SetHexaquarkPdg}, default 1600000); vertex
  Gaussian in $x,y$ (default $\sigma=4$~mm) clipped to \code{SetTrapRadius} (default 10~mm),
  uniform in $z$ over $\pm$\code{SetVertexZHalfRange} (default 22.5~mm). Vertex setters are in mm. \\
\code{AtPiGunGenerator} & \code{FairGenerator}; 1--3 pions (\code{SetMultiplicityRange}) with
  isotropic direction, momentum uniform in \code{SetMomentumRange} (default 0.10--0.30~GeV/c),
  random sign unless \code{SetCharge}; same vertex model as \code{AtPUMAGenerator}. \\
\bottomrule
\end{tabular}
\end{center}

\code{AtPUMAGenerator} and \code{AtPiGunGenerator} derive directly from \code{FairGenerator} and do
not take part in the beam/reaction alternation; every FairRoot event is a physics event. The
macro \url{macro/Simulation/ATTPC/PUMA/PiGun_sim.C} calls
\begin{verbatim}
AtVertexPropagator::Instance()->SetRndELoss(1e30);
\end{verbatim}
so that \code{AtTpc} never triggers a reaction on track 0.

%=====================================================================================
\section{Digitization}
\label{sec:sim:digi}

The digitization macro is a \code{FairRunAna} job with the transport output as
\code{FairFileSource} and an ASCII parameter file read into the runtime database. All digitization
tasks read the container \code{AtDigiPar}. The minimal chain is

\begin{verbatim}
auto mapping = std::make_shared<AtTpcMap>();
mapping->ParseXMLMap(mapParFile.Data());
mapping->GeneratePadPlane();

AtClusterizeTask *clusterizer = new AtClusterizeTask();
clusterizer->SetPersistence(kFALSE);

auto pulseModel = std::make_shared<AtPulse>(mapping);
AtPulseTask *pulse = new AtPulseTask(pulseModel);
pulse->SetPersistence(kTRUE);
pulse->SetSaveMCInfo();
\end{verbatim}

\subsection{AtDigiPar fields used by the simulation}

Table~\ref{tab:sim:digipar} lists the fields of \code{AtDigiPar} read by the digitization code.
All are mandatory in the parameter file (\code{AtDigiPar::getParams} aborts if one is missing)
except \code{ReverseDrift}, which is optional and defaults to 0.

\begin{table}[htbp]
\centering
\begin{tabular}{p{3.0cm}p{2.3cm}p{9.4cm}}
\toprule
Field & Unit & Used by \\
\midrule
\code{EIonize} & eV & \code{AtClusterize}: mean energy per ionization electron \\
\code{Fano} & -- & \code{AtClusterize}: width of the electron-number distribution \\
\code{DriftVelocity} & cm/$\mu$s & \code{AtClusterize} (drift time), \code{AtPulse}
  (pad-plane time bucket), space-charge models \\
\code{CoefT}, \code{CoefL} & cm$^2$/$\mu$s & \code{AtClusterize}: transverse and longitudinal
  diffusion \\
\code{ZPadPlane} & mm & \code{AtClusterize} (drift length), \code{AtPulse} (pad-plane time
  bucket), \code{AtLineChargeModel} \\
\code{ReverseDrift} & 0/1 & \code{AtClusterize}: beam enters through the pad plane \\
\code{Gain} & -- & \code{AtPulse}: mean micromegas gain (Polya) \\
\code{GETGain} & fC & \code{AtPulse}: full-scale range of the GET electronics \\
\code{PeakingTime} & ns & \code{AtPulse}: shaping time of the response function \\
\code{SamplingRate} & MHz code & \code{GetTBTime()}: 3, 6, 12, 25, 50, 100 map to 320, 160, 80,
  40, 20, 10~ns per time bucket; any other value returns $-1$ \\
\code{TBEntrance} & TB & \code{AtPulse}: time bucket of the detector entrance \\
\code{EField} & V/m & space-charge models \\
\bottomrule
\end{tabular}
\caption{AtDigiPar fields read by the simulation tasks.}
\label{tab:sim:digipar}
\end{table}

The number of time buckets is not read from the parameter file: \code{AtPulse::fNumTbs} is fixed
at 512 in the header.

\subsection{AtClusterizeTask and AtClusterize}

\code{AtClusterizeTask} reads the branch \code{AtTpcPoint} and writes \code{AtSimulatedPoint} (not
persistent by default; \code{SetPersistence(kTRUE)} keeps it). The physics is in the model class
passed to the constructor, \code{AtClusterize} by default; \code{SetClusterizeMethod} replaces it.

For each \code{AtMCPoint} in the volume named exactly \code{drift\_volume}, \code{AtClusterize}:

\begin{enumerate}
  \item converts the position to $(x,\,y)$ in mm and a drift time
    $t = |z_d|/v_d$ in $\mu$s, with $z_d = \mathrm{ZPadPlane}/10 - z$ (cm) for the normal
    detector and $z_d = z$ for \code{ReverseDrift = 1};
  \item if the point starts a new track or has zero energy loss, records its position and returns
    nothing. The energy deposited in the first point of each track is therefore not converted;
  \item draws the number of electrons $n$ from a Gaussian of mean $\Delta E/W$ and width
    $\sqrt{F\,\Delta E/W}$, clamped at zero (the source comment explains that an unclamped negative
    value wrapped to $\sim10^{19}$ and exhausted memory);
  \item places the $n$ electrons at equal intervals on the segment from the previous point to the
    current point, and displaces each by diffusion: a radial offset drawn from a Gaussian of width
    $\sigma_T = 10\sqrt{2 D_T t}$~mm at a uniform azimuth, and a time offset with
    $\sigma_L = \sqrt{2 D_L t}/v_d$~$\mu$s. Both widths are evaluated at the drift time of the
    current point.
\end{enumerate}

Each output \code{AtSimulatedPoint} is one electron (charge 1) with position $(x, y, t)$ and the
index of its parent \code{AtMCPoint}.

If \code{ZPadPlane} is shorter than the geometry drift volume (for example 971.7312~mm in the
a1975 deuterium parameter file against the 100~cm tube of \code{ATTPC\_D300torr\_v2}), points beyond
the pad plane get a drift time from the absolute value of a negative distance, i.e.\ they are folded
back into the chamber. The code does not warn about this.

\paragraph{AtClusterizeLine.} \code{AtClusterizeLine} is passed to the constructor of
\code{AtClusterizeTask} as a \code{std::make\_shared} object (the class
\code{AtClusterizeLineTask} is deprecated). It does not generate individual electrons. It writes one
\code{AtSimulatedLine} per MC step with the start and end points, the total number of electrons and
the transverse and longitudinal diffusion widths, to be consumed by \code{AtPulseLine}.

\subsection{AtPulseTask and AtPulse}

\code{AtPulseTask} reads \code{AtSimulatedPoint}, runs the pulse model, and writes an
\code{AtRawEvent} to the branch \code{AtRawEvent} (persistent by default;
\code{SetOutputBranch} renames it). The model is the \code{AtPulse} object passed to the
constructor; \code{AtPulse} needs the pad map, which must already have its inhibited pads set.

\paragraph{Charge collection.} For every simulated electron \code{AtPulse::AssignElectronsToPad}
finds the pad under $(x,y)$ (\code{AtMap::GetPadNum}); electrons outside the pad plane are
skipped. The arrival time is $t + \mathrm{TB}_\mathrm{pad}\cdot\Delta t_\mathrm{TB}$ with
\[
\mathrm{TB}_\mathrm{pad} = \mathrm{TBEntrance} -
\frac{\mathrm{ZPadPlane}/10}{\Delta t_\mathrm{TB}\, v_d}.
\]
The electron is multiplied by a gain drawn from a Polya distribution of mean \code{Gain} and shape
parameter 1. With \code{UseFastGain(true)} (default) and more than 10 electrons in one simulated
point, the gain is instead a Gaussian of mean \code{Gain} and width $\sigma_G/\sqrt{n}$; with
\code{AtClusterize} each point carries one electron, so the Polya sampling is used per electron.
Pads inhibited as \code{AtMap::InhibitType::kTotal} receive no charge; pads inhibited as
\code{kLowGain} have the gain multiplied by the factor set with \code{SetLowGain(val)} (default 0,
i.e.\ such pads also receive nothing unless the factor is set).

\paragraph{Trace formation.} The collected charge is histogrammed per pad in 512 time buckets.
\code{AtPulse::FillPad} convolves each bucket with the response function
(the setter \code{SetDoConvolution(false)} skips this and writes zero traces) and \code{ApplyNoise} converts to ADC by multiplying
by
\[
g_\mathrm{GET} = \frac{1.602\times10^{-19}\times 4096}{\mathrm{GETGain}\times10^{-15}}
\quad\text{ADC per electron.}
\]
The default response function is \code{ElectronicResponse::AtNominalResponse} with the peaking time
from \code{AtDigiPar}:
$R(t) = e^{-3t/\tau}\,\sin(t/\tau)\,(t/\tau)^3$ (the code uses $2.718$ for $e$). Other response
functions are in \url{AtTools/ElectronicResponse/AtElectronicResponse.h}:
\code{AtReducedTimeResponse} (a user function of $t/\tau$), \code{AtVectorResponse} (a sampled trace)
and \code{AtFileResponse} (a trace read from a text file); any of them can be passed as the second
constructor argument of \code{AtPulse}. The response functions are not normalized; the comment in
the header notes that changing the response therefore changes the effective gain.
The class \code{AtVectorResponse} in \url{AtDigitization/} is deprecated in favour of
\code{AtFileResponse}.

With \code{SetSaveCharge(true)} (default) each pad also gets an augment named \code{"Q"} holding the
collected charge per time bucket. Pads are marked valid, pedestal-subtracted, and carry their pad
centre.

\paragraph{Setters of AtPulse.}

\begin{center}
\begin{tabular}{p{5.9cm}p{1.5cm}p{7.2cm}}
\toprule
Setter & Default & Effect \\
\midrule
\code{SetNoiseSigma(double)} & 0 & Gaussian baseline noise (ADC) added to every time bucket of
  every pad that has signal. \\
\code{SetNoiseAllPads(bool)} & false & If true and the noise sigma is non-zero, also writes
  pure-noise traces on every non-inhibited pad without signal. \\
\code{SetNoiseKeepThreshold(double)} & 0 & With \code{SetNoiseAllPads}, drop noise-only pads whose
  maximum is below this value (saves disk; no effect on hits if below the PSA threshold). \\
\code{SetADCCeiling(double)} & 0 (off) & Clip every sample, after noise, at this ADC value
  (GET saturation). \\
\code{SetLowGain(double)} & 0 & Gain factor for \code{kLowGain} pads. \\
\code{UseFastGain(bool)} & true & Gaussian approximation of the summed Polya gain for points with
  more than 10 electrons. \\
\code{SetSaveCharge(bool)} & true & Store the \code{"Q"} augment. \\
\code{SetDoConvolution(bool)} & true & Convolve with the response function. \\
\bottomrule
\end{tabular}
\end{center}

\paragraph{Setters of AtPulseTask.} \code{SetPersistence} (default true),
\code{SetOutputBranch} (default \code{"AtRawEvent"}), \code{SetPersistenceAtTpcPoint} (default
false: the \code{AtTpcPoint} branch is re-registered and, unless this is set, not copied to the
digitization output), and \code{SetSaveMCInfo()}. The last one fills a map
pad number $\to$ \code{AtMCPoint} index in the \code{AtRawEvent} (\code{SetSimMCPointMap}).
\code{AtPSAtask} finds the \code{AtTpcPoint} branch in the input and attaches to every hit an
\code{AtHit::MCSimPoint} with the point index, track ID, \code{EIni}, energy loss, \code{AIni}, $A$
and $Z$. This is how a simulation analysis matches hits to the generated particle.
As written, \code{AtPulseTask::saveMCInfo} compares the new point with itself when checking for
duplicates, so in practice each pad keeps only the first \code{AtMCPoint} that reached it; when two
tracks share a pad, only the first one is recorded.

\subsection{AtPulseLine and AtPulseGADGET}

\code{AtPulseLine} (an \code{AtPulse} subclass, used with \code{AtClusterizeLine}) integrates each
\code{AtSimulatedLine} instead of handling single electrons: the transverse distribution over pads is
sampled with \code{SetNumIntegrationPoints} random throws (default 1000), and the longitudinal
Gaussian is integrated analytically over time buckets within \code{SetNumSigmaToIntegrateZ} sigmas
(default 3). It uses the fast Gaussian gain. The old \code{AtPulseLineTask} is not compiled
(commented out in \url{AtDigitization/CMakeLists.txt}).

\code{AtPulseGADGET} is a pulse model for the GADGET detector. Each electron's charge is dispersed
over a grid of neighbouring pads with an RC charge-dispersion model; the grid spacing (2.2~mm) and
the RC constants are hard-coded members. \code{SetAdjecentPads} changes the pad offsets, but the
grid size \code{Items} is computed once from the default of one adjacent pad, so the grid stays
$3\times3$.

\subsection{Space charge and field distortion}

\code{AtSpaceChargeTask} exists in \url{AtDigitization/} but is not in the library
(\code{CMakeLists.txt} does not list it) and its \code{processPoint} is empty. Space-charge
distortion in the simulation is available only in \code{AtSimpleSimulation} through
\code{SetSpaceChargeModel}. The models derive from \code{AtSpaceChargeModel}
(\url{AtTools/}) with \code{ApplySpaceCharge} (simulation) and \code{CorrectSpaceCharge} (data):

\begin{itemize}
  \item \code{AtLineChargeModel}: line charge \code{SetLambda} (default $1.54\times10^{-8}$~C/m)
    along the beam, drift field \code{SetDriftField} (default 70000~V/m), beam radius
    \code{SetBeamRadius} (default 0.02~m, no correction inside); \code{LoadParameters} takes the
    field from \code{EField} and the length from \code{ZPadPlane}.
  \item \code{AtRadialChargeModel}: arbitrary radial field (a callable of $\rho$, $z$ in cm returning
    V/cm, e.g.\ \code{AtLineChargeZDep}); the electron path is integrated with step
    \code{SetStepSize} (default $10^{-4}$~$\mu$s); \code{LoadParameters} reads \code{EField} and
    \code{DriftVelocity}.
  \item \code{AtEDistortionModel}: look-up tables of $z$, radial and transverse corrections
    set with \code{SetCorrectionMaps}; its \code{ApplySpaceCharge} returns an empty point, so it can only
    be used for correction.
\end{itemize}

\subsection{AtTriggerTask}

\code{AtTriggerTask} emulates the GET multiplicity trigger on already digitized and
pulse-analysed events. It reads \code{AtEventH} and \code{AtRawEvent} (so it must run after the PSA)
and copies events that fire the trigger into the branches \code{Accepted\_AtEventH} and
\code{Accepted\_AtRawEvent}. Its parameters come from the container \code{AtTriggerPar}
(\url{parameters/AT.trigger.par}): \code{Write\_clock}, \code{Read\_clock}, \code{Master\_clock},
\code{Pad\_thresh\_MSB}, \code{Pad\_thresh\_LSB}, \code{Trigger\_signal\_width},
\code{Trigger\_discriminator\_fraction}, \code{Multiplicity\_threshold},
\code{Multiplicity\_window}, \code{Trigger\_height}. The pad map is given with
\code{SetAtMap(path)}. Two points visible in the code: \code{Init} passes the write and read clocks
to \code{AtTrigger::SetTriggerParameters} in the order \code{(write, read, ...)} while that function
declares \code{(read, write, ...)}; and the accepted arrays are not cleared between events (the
\code{Delete()} calls are commented out). No macro in the repository uses this task.

\subsection{AtSimpleSimulation}
\label{sec:sim:simple}

\code{AtSimpleSimulation} (units MeV, mm, MeV/c) steps a particle in straight lines of
\code{SetDistanceStep} (default 1~mm), takes the energy loss from an \code{AtTools::AtELossModel}
registered per $(Z,A)$ with \code{AddModel(Z, A, model)}, and writes one \code{AtMCPoint} per step
to the branch registered by \code{RegisterBranch} (default \code{"AtTpcPoint"}). Direction is not
changed by the energy loss and no magnetic field is applied. With a space-charge model set, each
point is displaced by \code{ApplySpaceCharge} after converting $z$ to the analysis frame
($1000-z$, hard-coded). The FairTask wrapper is \code{AtTestSimulation}. The output feeds the same
\code{AtClusterizeTask}/\code{AtPulseTask} chain.

%=====================================================================================
\section{Worked example: 16C(d,t)15C}
\label{sec:sim:example}

The a1975 deuterium-target simulation in \url{macro/Simulation/ATTPC/16C_dt/} measures the
acceptance of the triton channel for five states of 15C. It is split into generation
(\code{C16\_dt\_sim.C}, driven by \code{run\_states\_dt.sh}) and digitization plus reconstruction
(\code{run\_reco\_C16dt.C}, driven by \code{run\_reco\_batch\_dt.sh}).

\subsection{Generation: C16\_dt\_sim.C}

Signature and conditions, as stated in the macro header:

\begin{verbatim}
void C16_dt_sim(Int_t nEvents = 2000, Double_t thetaMinCM = 2.0,
                Double_t thetaMaxCM = 178.0, TString mcEngine = "TGeant4",
                Double_t bFieldkG = -28.5,
                TString outFile = "./data/attpcsim_dt.root",
                Double_t resEx = 0.0, UInt_t seed = 0, Double_t Ebeam = 184.17)
//  B = 2.85 T, D2 at 300 torr, beam 16C at 184.17 MeV = 11.5 MeV/u
\end{verbatim}

The seed is applied first:

\begin{verbatim}
if (seed != 0)
   gRandom->SetSeed(seed);
std::cout << "RNG seed requested: " << seed << std::endl;
\end{verbatim}

The beam momentum is computed from the kinetic energy and passed per nucleon:

\begin{verbatim}
const Double_t u = 931.49401, mBeamMeV = 16.0147 * u;
// GeV/c
Double_t pTot = std::sqrt(Ebeam * Ebeam + 2 * Ebeam * mBeamMeV) / 1000.0;
Double_t pz = pTot / a; // GeV/c per nucleon
Double_t Bmass = 16.0147;   // amu (repo convention)
Double_t maxELoss = 10.0;   // MeV
AtTPCIonGenerator *ionGen =
   new AtTPCIonGenerator("Ion", z, a, q, m, px, py, pz, BExcEner, Bmass,
                         NomEnergy, maxELoss);
ionGen->SetSpotRadius(0, -100, 0);
\end{verbatim}

The beam starts on axis at $z=-100$~cm. The two-body generator gets the four entries
beam (16C), target (d, 2.0141017781~u), heavy residual (15C, 15.0105993~u, excitation
\code{resEx}) and light ejectile (t, 3.0160492779~u):

\begin{verbatim}
AtTPC2Body *TwoBody =
   new AtTPC2Body("TwoBody", &Zp, &Ap, &Qp, mult, &Pxp, &Pyp, &Pzp, &Mass,
                  &ExE, ResEner, ThetaMinCMS, ThetaMaxCMS);
primGen->AddGenerator(TwoBody);
run->SetGenerator(primGen);
run->SetStoreTraj(kTRUE);
\end{verbatim}

The macro notes that states of 15C above the neutron separation energy (1.218~MeV) are transported
as stable ions; the simulation describes the triton but must not be used for the recoil.

\paragraph{run\_states\_dt.sh.} The driver generates the five states (g.s., 0.740, 3.103, 4.657,
6.358~MeV) with one seed, the residual centre-of-mass range 2--70$^\circ$, and the data field sign:

\begin{verbatim}
root -b -q -l "C16_dt_sim.C($NEV,2.,70.,\"TGeant4\",-28.5,\
\"$OUT/${J}_sim.root\",$EX,$SEED)" > "$OUT/${J}_gen.log" 2>&1
\end{verbatim}

A job is marked complete only if the output is non-empty, the log contains
\code{"RNG seed requested: \$SEED"} and \code{"Macro finished successfully"}; otherwise no
\code{.marker} file is written and the digitization driver skips that state.

\subsection{Digitization and reconstruction: run\_reco\_C16dt.C}

The macro reads the data's own parameter file (or a gain variant of it), sets the beam hole,
digitizes, and runs the same PSA, cleaning and pattern-recognition settings as the data
production:

\begin{verbatim}
#include "../AtBeamHole.h"
...
auto mapping = std::make_shared<AtTpcMap>();
mapping->ParseXMLMap(mapParFile.Data());
mapping->GeneratePadPlane();
AtSim::InhibitBeamHole(mapping, holeR);          // holeR = 20 mm for a1975

AtClusterizeTask *clusterizer = new AtClusterizeTask();
clusterizer->SetPersistence(kFALSE);

auto pulseModel = std::make_shared<AtPulse>(mapping);
if (adcCeiling > 0)
   pulseModel->SetADCCeiling(adcCeiling);
AtPulseTask *pulse = new AtPulseTask(pulseModel);
pulse->SetPersistence(kTRUE);
pulse->SetSaveMCInfo();

auto p = new AtPSAMultiFit();
p->SetThreshold(thr);
p->SetPeakingTime(0.720);
...
fRun->AddTask(clusterizer);
fRun->AddTask(pulse);
fRun->AddTask(psaTask);
fRun->AddTask(cleanTask);
fRun->AddTask(praTask);
\end{verbatim}

The parameter file defaults to \url{parameters/ATTPC.a1975_deuterium_dv1104.par}: drift velocity
1.10424~cm/$\mu$s, \code{TBEntrance} 560, \code{ZPadPlane} 971.7312~mm, \code{SamplingRate} 6
(160~ns), \code{EIonize} 15.603~eV, \code{Fano} 0.24, \code{CoefL} = \code{CoefT} =
0.0009~cm$^2$/$\mu$s, \code{GETGain} 1000~fC, \code{PeakingTime} 720~ns. The macro header explains
that drift velocity, \code{TBEntrance} and \code{ZPadPlane} must equal the data values, otherwise the
$z$ scale of the simulation is not that of the data.

\paragraph{run\_reco\_batch\_dt.sh and the gain.} \code{AtDigiPar} has no gain setter, so the gain is
scanned through parameter files: \code{make\_gain\_pars.sh} writes
\url{ATTPC.a1975_deuterium_dv1104_g<N>.par} files that differ from the production file only in the
\code{Gain} line. The batch driver uses \code{Gain = 35000} by default and runs the five states in
parallel:

\begin{verbatim}
./run_reco_batch_dt.sh [nparallel] [gain] [simDir]
root -b -l -q "run_reco_C16dt.C(\"$in\",\"$out\",\"$PAR\")"
\end{verbatim}

Its header gives the reason for 35000, a scan of the median number of hits per non-beam track on 200
events of the ground-state sample compared with data run 0016 (median 76):

\begin{center}
\begin{tabular}{p{2.0cm}p{3.6cm}p{2.0cm}}
\toprule
Gain & median hits/track & sim/data \\
\midrule
10000 & 39 & 0.51 \\
25000 & 62 & 0.82 \\
35000 & 75 & 0.99 \\
50000 & 90 & 1.18 \\
150000 & 118 & 1.55 \\
400000 & 135 & 1.78 \\
\bottomrule
\end{tabular}
\end{center}

The header also records that the event efficiency saturates above about 25000, so the gain chooses
how much of each track survives rather than how many events do, and that the scan is not
PID-selected; it rejects 10000 and 150000 but does not fix 35000 to two digits.

%=====================================================================================
\section{Known pitfalls}
\label{sec:sim:pitfalls}

\begin{enumerate}
  \item \textbf{maxELoss sets the reaction rate and the vertex distribution.} The threshold is
    uniform in accumulated energy loss of the beam in the gas, in $[0,\mathrm{maxELoss}]$. A value
    larger than the beam loses over the chamber leaves many beam events without a reaction (empty
    reaction events); a value much smaller piles all vertices near the window. If the constructor
    argument is omitted, the nominal beam energy is used. Check the vertex-$z$ distribution in the
    MC truth (\url{macro/Simulation/ATTPC/16C_dt/dt_truth_check.C} does this for the (d,t)
    sample).
  \item \textbf{Two FairRoot events per reaction.} Even entries are beam events and odd entries are
    reaction events; \code{nEvents} counts both. Beam and products are never in the same event.
  \item \textbf{Mass units.} The macros pass the ion mass in amu to \code{AtTPCIonGenerator} and
    \code{AtTPC2Body}, while \code{FairIon} documents the mass argument in GeV. The AT-TPC code
    that reads the stored mass treats it as amu. Whatever the transport engine does with it, the
    beam energy must be verified from the MC truth (the \code{MCTrack} four-momentum of the
    primary) rather than assumed; \url{macro/Simulation/ATTPC/14C_pp/check_beam_C14.C} reads it
    back. \url{macro/Simulation/ATTPC/17C_dp/RESULTS.md} reports the 17C beam momentum and kinetic
    energy recovered exactly with this convention.
  \item \textbf{Seeds for parallel jobs.} Without \code{gRandom->SetSeed(seed)} parallel jobs
    produce identical events. \code{C16\_dt\_sim.C} prints the requested seed because, as its
    comment states, \code{gRandom->GetSeed()} on \code{TRandom3} does not return the seed.
    Give every job a distinct seed.
  \item \textbf{Field sign.} Generate with the field sign of the data
    (\code{bFieldkG = -28.5} for a1975), so that the sense of rotation in the simulation matches
    the data. The comment in \code{C16\_dt\_sim.C} warns that otherwise any step that infers the
    direction from the curvature rejects most simulated tracks.
  \item \textbf{z frame.} The PSA writes hit $z$ measured from the pad plane; the fitters convert
    with $z_\mathrm{lab} = \mathrm{ZPadPlane} - z_\mathrm{digi}$ (the comments in
    \url{AtReconstruction/AtPulseAnalyzer/AtPSA.cxx} and \url{AtParameter/AtDigiPar.h}). With
    \code{ZPadPlane} = 1000~mm this is $z_\mathrm{lab} = 1000 - z_\mathrm{digi}$. The header of
    \code{run\_reco\_C16dt.C} states that a fit of the simulation output needs \code{bFieldSign = +1}
    where the data uses $-1$, and that the reconstructed polar angle is $180^\circ - \theta_\mathrm{true}$;
    account for this when matching to truth.
  \item \textbf{Gain must be matched to the data.} The digitized gain sets how many pads cross the
    PSA threshold and so the number of hits per track that pattern recognition and PID see. Match
    hits per track to the data before trusting an acceptance (see the scan above); do not inherit a
    gain from another channel's simulation parameter file.
  \item \textbf{Truncated heavy residual in old productions.} Before commit \code{ac1b8fb1}
    the reaction threshold fired a second time on the heavy residual (track ID 0 in the reaction
    event), stopping it after roughly the vertex depth. The light ejectile was not affected. Files
    produced before the fix carry half-length residuals.
  \item \textbf{No saturation by default.} \code{AtPulse} does not saturate unless
    \code{SetADCCeiling} is set; simulated high-\dedx{} tracks then keep charge that the GET
    electronics would clip. The commit that added the ceiling (\code{f5d5031f}) states that
    \code{AtPSAMultiFit} maps a clipped pulse to about 1.34 times the clip value, and that the data's
    saturation shape still has to be measured; \code{run\_reco\_C16dt.C} leaves it off by default.
  \item \textbf{Beam hole.} The real pad plane does not read out the beam region, but the
    simulation digitizes the beam in full. \url{macro/Simulation/ATTPC/AtBeamHole.h} provides
    \code{AtSim::InhibitBeamHole(map, holeR)}, which inhibits as \code{kTotal} every pad whose centre
    is within \code{holeR} (mm) of the axis; call it after \code{GeneratePadPlane()} and before
    building \code{AtPulse}. Its header reports that on a1954 14C(p,p') without the hole 67\% of the
    recoil protons were merged into the beam cluster, and that the radius should be measured per
    experiment from the hit-radius distribution (30~mm for a1954, 20~mm for a1975 in
    \code{run\_reco\_C16dt.C}).
  \item \textbf{ZPadPlane versus geometry length.} Points beyond \code{ZPadPlane} are folded back by
    the absolute value in \code{AtClusterize}; keep the parameter and the geometry consistent or
    expect a small amount of charge at mirrored drift times.
  \item \textbf{One MC point per pad.} With \code{SetSaveMCInfo()} only the first \code{AtMCPoint}
    reaching a pad is recorded (see the description of \code{saveMCInfo}). Hit-to-track truth
    matching is unreliable where tracks overlap.
\end{enumerate}
